\documentclass[12pt,a4paper]{article}

\input{/home/tabaka/Documents/GitHub/Defs/defs.tex}

\setcounter{zadaniaRozdzial}{3}
\setcounter{section}{0}

\begin{document}
	
	\begin{center}
		\Large Wymagania edukacyjne z matematyki w klasie 1 Liceum Ogólnokształcącego na poszczególne oceny.
	\end{center}
	
	Wymagania edukacyjne na poszczególne oceny są stopniowane.
	Uzyskanie danej oceny wymaga spełnienia wymagań przewidzianych
	dla wszystkich ocen niższych, z uwzględnieniem wymagań
	określonych dla danej oceny.
	
	\section{Liczby rzeczywiste}
	
	\begin{longtable}{|p{16cm}|}
		\hline
		\textbf{Wymagania edukacyjne na ocenę dopuszczającą} \\ \hline
		
		\endfirsthead
		
		%\textbf{Wymagania edukacyjne na ocenę dopuszczającą -- ciąg dalszy} \\ \hline
		\endhead
		
		\textbullet\ podaje przykłady liczb naturalnych, całkowitych, wymiernych i niewymiernych oraz przyporządkowuje liczbę do odpowiedniego zbioru liczb; \\
		
		\textbullet\ rozróżnia liczby pierwsze i liczby złożone; \\
		
		\textbullet\ wskazuje liczby podzielne np. przez $2$, $3$, $4$, $5$, $9$, $10$; \\
		
		\textbullet\ podaje dzielniki danej liczby naturalnej; \\
		
		\textbullet\ przedstawia liczby naturalne w postaci iloczynu liczb pierwszych; \\
		
		\textbullet\ podaje liczbę przeciwną oraz odwrotną do danej liczby; \\
		
		\textbullet\ porównuje liczby wymierne; \\
		
		\textbullet\ podaje przykład liczby wymiernej zawartej między dwiema danymi liczbami wymiernymi; \\
		
		\textbullet\ zaznacza na osi liczbowej daną liczbę wymierną oraz odczytuje z osi liczbowej współrzędne danego punktu; \\
		
		\textbullet\ przedstawia liczby wymierne w różnych postaciach; \\
		
		\textbullet\ wyznacza przybliżenia dziesiętne danej liczby rzeczywistej z zadaną dokładnością, również przy użyciu kalkulatora, oraz określa, czy dane przybliżenie jest przybliżeniem z nadmiarem, czy z niedomiarem; \\
		
		\textbullet\ wyznacza rozwinięcia dziesiętne ułamków zwykłych oraz zamienia skończone rozwinięcia dziesiętne na ułamki zwykłe; \\
		
		\textbullet\ wykonuje proste działania w zbiorze liczb wymiernych; \\
		
		\textbullet\ oblicza wartość pierwiastka dowolnego stopnia z liczby nieujemnej oraz wartość pierwiastka nieparzystego stopnia z liczby rzeczywistej w prostych przypadkach; \\
		
		\textbullet\ wyłącza czynnik przed znak pierwiastka kwadratowego oraz włącza czynnik pod znak pierwiastka kwadratowego w prostych przypadkach; \\
		
		\textbullet\ wykonuje działania na pierwiastkach tego samego stopnia, stosując odpowiednie twierdzenia, w prostych przypadkach; \\
		
		\textbullet\ usuwa niewymierność z mianownika wyrażenia typu $\frac{1}{\sqrt{a}}$; \\
		
		\textbullet\ oblicza potęgi o wykładnikach wymiernych; \\
		
		\textbullet\ zapisuje daną liczbę w postaci potęgi o wykładniku całkowitym; \\
		
		\textbullet\ zapisuje daną liczbę w postaci potęgi o danej podstawie w prostych przypadkach; \\
		
		\textbullet\ oblicza logarytm liczby w prostych przypadkach; \\
		
		\textbullet\ oblicza procent danej liczby; \\
		
		\textbullet\ oblicza, jakim procentem jednej liczby jest druga liczba; \\
		
		\textbullet\ wyznacza liczbę, gdy dany jest jej procent. \\ \hline
		
		\textbf{Wymagania edukacyjne na ocenę dostateczną} \\ \hline
		
		\textbullet\ oblicza największy wspólny dzielnik (NWD) i najmniejszą wspólną wielokrotność (NWW) danych liczb naturalnych; \\
		
		\textbullet\ porównuje liczby niewymierne; \\
		
		\textbullet\ podaje przykład liczby niewymiernej zawartej między dwiema danymi liczbami; \\
		
		\textbullet\ zamienia ułamki dziesiętne okresowe, np. $0,(2)$ i $0,(02)$, na ułamki zwykłe; \\
		
		\textbullet\ wykonuje działania łączne w zbiorze liczb rzeczywistych; \\
		
		\textbullet\ oblicza wartość pierwiastka dowolnego stopnia z liczby nieujemnej oraz wartość pierwiastka nieparzystego stopnia z liczby rzeczywistej; \\
		
		\textbullet\ wykonuje działania na pierwiastkach tego samego stopnia, stosując odpowiednie twierdzenia; \\
		
		\textbullet\ przekształca wyrażenia zawierające pierwiastki kwadratowe i oblicza ich wartości; \\
		
		\textbullet\ konstruuje odcinki o długościach niewymiernych, np. $\sqrt{5}$; \\
		
		\textbullet\ zapisuje daną liczbę w postaci potęgi o wykładniku wymiernym; \\
		
		\textbullet\ zapisuje daną liczbę w postaci potęgi o danej podstawie; \\
		
		\textbullet\ upraszcza wyrażenia, stosując prawa działań na potęgach; \\
		
		\textbullet\ porównuje liczby przedstawione w postaci potęg; \\
		
		\textbullet\ stosuje równości wynikające z definicji logarytmu do prostych obliczeń; \\
		
		\textbullet\ stosuje twierdzenia o logarytmie iloczynu, ilorazu i potęgi w prostych przypadkach; \\
		
		\textbullet\ zwiększa i zmniejsza liczbę o dany procent; \\
		
		\textbullet\ oblicza, o ile procent jedna liczba jest większa lub mniejsza od drugiej; \\
		
		\textbullet\ posługuje się procentami w rozwiązywaniu prostych zadań praktycznych. \\ \hline
		
		\textbf{Wymagania edukacyjne na ocenę dobrą} \\ \hline
		
		\textbullet\ stosuje ogólny zapis liczb naturalnych parzystych, nieparzystych, podzielnych przez $3$ itp.; \\
		
		\textbullet\ przeprowadza dowody twierdzeń dotyczących podzielności liczb w prostych przypadkach; \\
		
		\textbullet\ wykorzystuje dzielenie z resztą do przedstawiania liczby naturalnej w postaci $a \cdot k + r$; \\
		
		\textbullet\ wykonuje działania łączne na liczbach rzeczywistych w trudniejszych przypadkach; \\
		
		\textbullet\ zamienia ułamek dziesiętny okresowy na ułamek zwykły; \\
		
		\textbullet\ wyznacza wskazaną cyfrę po przecinku w rozwinięciu dziesiętnym okresowym danej liczby w prostych przypadkach; \\
		
		\textbullet\ porównuje pierwiastki bez użycia kalkulatora; \\
		
		\textbullet\ oblicza wartości wyrażeń arytmetycznych zawierających pierwiastki, stosując prawa działań na pierwiastkach; \\
		
		\textbullet\ wyłącza czynnik przed znak pierwiastka dowolnego stopnia oraz włącza czynnik pod znak pierwiastka dowolnego stopnia; \\
		
		\textbullet\ konstruuje odcinki o długościach niewymiernych, np. $\sqrt{15}$; \\
		
		\textbullet\ stosuje działania na pierwiastkach do obliczania pól czworokątów; \\
		
		\textbullet\ usuwa niewymierność z mianownika wyrażenia typu $\frac{1}{\sqrt[3]{a}}$; \\
		
		\textbullet\ upraszcza wyrażenia, stosując prawa działań na potęgach w trudniejszych przypadkach; \\
		
		\textbullet\ porównuje liczby przedstawione w postaci potęg w trudniejszych przypadkach; \\
		
		\textbullet\ stosuje twierdzenia o logarytmie iloczynu, ilorazu i potęgi do obliczeń; \\
		
		\textbullet\ rozwiązuje złożone zadania tekstowe, wykorzystując obliczenia procentowe. \\ \hline
		
		\textbf{Wymagania edukacyjne na ocenę bardzo dobrą} \\ \hline
		
		\textbullet\ przeprowadza dowody twierdzeń dotyczących podzielności liczb i reszt z dzielenia w trudniejszych przypadkach; \\
		
		\textbullet\ wyznacza wskazaną cyfrę po przecinku w rozwinięciu dziesiętnym okresowym danej liczby; \\
		
		\textbullet\ przeprowadza dowody twierdzeń dotyczących logarytmu iloczynu, ilorazu i potęgi; \\
		
		\textbullet\ stosuje twierdzenia dotyczące logarytmu iloczynu, ilorazu i potęgi do wykazywania równości wyrażeń. \\ \hline
		
		\textbf{Wymagania edukacyjne na ocenę celującą} \\ \hline
		
		\textbullet\ rozwiązuje zadania o znacznym stopniu trudności dotyczące liczb rzeczywistych. \\ \hline
		
	\end{longtable}
	
	\section{Język matematyki}
	
	\begin{longtable}{|p{16cm}|}
		\hline
		\textbf{Wymagania edukacyjne na ocenę dopuszczającą} \\ \hline
		
		\textbullet\ posługuje się pojęciami zbioru, podzbioru, zbioru skończonego i zbioru nieskończonego; \\
		
		\textbullet\ wymienia elementy danego zbioru; \\
		
		\textbullet\ posługuje się pojęciami sumy i iloczynu zbiorów; \\
		
		\textbullet\ zaznacza na osi liczbowej przedziały liczbowe; \\
		
		\textbullet\ wyznacza sumę i iloczyn przedziałów liczbowych oraz zaznacza je na osi liczbowej w prostych przypadkach; \\
		
		\textbullet\ rozwiązuje proste nierówności liniowe i sprawdza, czy dana liczba spełnia daną nierówność; \\
		
		\textbullet\ zaznacza na osi liczbowej zbiór rozwiązań nierówności liniowej; \\
		
		\textbullet\ wyłącza wskazany jednomian przed nawias w sumie algebraicznej; \\
		
		\textbullet\ zapisuje związki między wielkościami za pomocą wyrażeń algebraicznych w prostych przypadkach; \\
		
		\textbullet\ stosuje wzory skróconego mnożenia do wyznaczania kwadratu sumy lub różnicy oraz różnicy kwadratów; \\
		
		\textbullet\ oblicza wartość bezwzględną liczby rzeczywistej. \\ \hline
		
		\textbf{Wymagania edukacyjne na ocenę dostateczną} \\ \hline
		
		\textbullet\ opisuje symbolicznie dane zbiory w prostych przypadkach; \\
		
		\textbullet\ posługuje się pojęciem różnicy zbiorów; \\
		
		\textbullet\ wyznacza różnicę przedziałów liczbowych i zaznacza ją na osi liczbowej; \\
		
		\textbullet\ rozwiązuje nierówności liniowe i zaznacza ich rozwiązania na osi liczbowej; \\
		
		\textbullet\ zapisuje zbiory w postaci przedziałów liczbowych, np. $A=\{x\in\mathbb{R}:x\geq -4\land x<1\}=[-4;1)$; \\
		
		\textbullet\ mnoży sumy algebraiczne przez siebie i redukuje wyrazy podobne w otrzymanej sumie; \\
		
		\textbullet\ stosuje wzory skróconego mnożenia do przekształcania wyrażeń algebraicznych w prostych przypadkach; \\
		
		\textbullet\ usuwa niewymierność z mianownika ułamka, gdy mianownik ma postać $a\sqrt{b}$; \\
		
		\textbullet\ stosuje przekształcenia wyrażeń algebraicznych do rozwiązywania prostych równań i nierówności; \\
		
		\textbullet\ stosuje geometryczną interpretację wartości bezwzględnej do rozwiązywania prostych równań i nierówności typu $|x|=a, \: |x|<a$. \\ \hline
		
		\textbf{Wymagania edukacyjne na ocenę dobrą} \\ \hline
		
		\textbullet\ zaznacza na osi liczbowej zbiory liczb spełniających układ nierówności liniowych z jedną niewiadomą; \\
		
		\textbullet\ zapisuje związki między wielkościami za pomocą wyrażeń algebraicznych; \\
		
		\textbullet\ przeprowadza proste dowody, stosując działania na wyrażeniach algebraicznych; \\
		
		\textbullet\ stosuje wzory skróconego mnożenia do przekształcania wyrażeń algebraicznych; \\
		
		\textbullet\ stosuje wzory skróconego mnożenia do wykonywania działań na liczbach zapisanych w postaci wyrażeń algebraicznych; \\
		
		\textbullet\ usuwa niewymierność z mianownika wyrażenia typu $\frac{a}{\sqrt{b}\pm c\sqrt{d}}$; \\
		
		\textbullet\ stosuje przekształcenia algebraiczne do rozwiązywania równań i nierówności; \\
		
		\textbullet\ stosuje nierówności pierwszego stopnia z jedną niewiadomą do rozwiązywania zadań osadzonych w kontekście praktycznym; \\
		
		\textbullet\ upraszcza wyrażenia zawierające wartość bezwzględną, w tym stosuje własność $\sqrt{x^2}=|x|$; \\
		
		\textbullet\ stosuje geometryczną interpretację wartości bezwzględnej do rozwiązywania równań typu $|x-a|=b$; \\
		
		\textbullet\ wyprowadza wzory skróconego mnożenia. \\ \hline
		
		\textbf{Wymagania edukacyjne na ocenę bardzo dobrą} \\ \hline
		
		\textbullet\ wyznacza dopełnienie zbioru; \\
		
		\textbullet\ wykonuje złożone działania na przedziałach liczbowych; \\
		
		\textbullet\ stosuje wzory skróconego mnożenia do dowodzenia twierdzeń; \\
		
		\textbullet\ stosuje przekształcenia algebraiczne do rozwiązywania równań i nierówności w trudniejszych przypadkach; \\
		
		\textbullet\ stosuje nierówności pierwszego stopnia z jedną niewiadomą do rozwiązywania zadań osadzonych w kontekście praktycznym w trudniejszych przypadkach; \\
		
		\textbullet\ upraszcza wyrażenia zawierające wartość bezwzględną w trudniejszych przypadkach. \\ \hline
		
		\textbf{Wymagania edukacyjne na ocenę celującą} \\ \hline
		
		\textbullet\ rozwiązuje zadania o znacznym stopniu trudności dotyczące zbiorów, przekształcania wyrażeń algebraicznych i własności wartości bezwzględnej. \\ \hline
		
	\end{longtable}
	
	\section{Układy równań}
	
	\begin{longtable}{|p{16cm}|}
		\hline
		
		\textbf{Wymagania edukacyjne na ocenę dopuszczającą} \\ \hline
		
		\textbullet\ podaje przykładowe rozwiązania równania liniowego z dwiema niewiadomymi; \\
		
		\textbullet\ sprawdza, czy dana para liczb spełnia dany układ równań; \\
		
		\textbullet\ wyznacza wskazaną zmienną z równania liniowego z dwiema niewiadomymi; \\
		
		\textbullet\ rozwiązuje układy równań metodą podstawiania w prostych przypadkach, gdy równania są odpowiednio uporządkowane; \\
		
		\textbullet\ rozwiązuje układy równań metodą przeciwnych współczynników w prostych przypadkach, gdy równania są odpowiednio uporządkowane; \\
		
		\textbullet\ rozpoznaje układ oznaczony, nieoznaczony i sprzeczny. \\ \hline
		
		\textbf{Wymagania edukacyjne na ocenę dostateczną} \\ \hline
		
		\textbullet\ dopisuje do danego równania drugie równanie tak, aby wskazana para liczb była rozwiązaniem układu; \\
		
		\textbullet\ rozwiązuje układy równań metodą podstawiania; \\
		
		\textbullet\ rozwiązuje układy równań metodą przeciwnych współczynników; \\
		
		\textbullet\ określa, czy dany układ równań jest oznaczony, nieoznaczony czy sprzeczny; \\
		
		\textbullet\ stosuje układy równań liniowych do rozwiązywania prostych zadań tekstowych. \\ \hline
		
		\textbf{Wymagania edukacyjne na ocenę dobrą} \\ \hline
		
		\textbullet\ zapisuje w postaci układu równań informacje przedstawione w zadaniu tekstowym; \\
		
		\textbullet\ dobiera współczynniki liczbowe w układzie równań tak, aby wskazana para liczb była jego rozwiązaniem; \\
		
		\textbullet\ dopisuje drugie równanie tak, aby otrzymany układ był sprzeczny, oznaczony lub nieoznaczony; \\
		
		\textbullet\ rozwiązuje układy równań w trudniejszych przypadkach, stosując przekształcenia algebraiczne i wzory skróconego mnożenia. \\ \hline
		
		\textbf{Wymagania edukacyjne na ocenę bardzo dobrą} \\ \hline
		
		\textbullet\ przedstawia zbiór rozwiązań układu nieoznaczonego; \\
		
		\textbullet\ stosuje układy równań do rozwiązywania złożonych zadań tekstowych, w tym dotyczących prędkości oraz wielkości wyrażonych za pomocą procentów, takich jak stężenia roztworów i lokaty bankowe. \\ \hline
		
		\textbf{Wymagania edukacyjne na ocenę celującą} \\ \hline
		
		\textbullet\ rozwiązuje zadania o znacznym stopniu trudności dotyczące układów równań, w tym układy równań z trzema niewiadomymi; \\
		
		\textbullet\ stosuje układy równań do rozwiązywania trudniejszych zadań tekstowych. \\ \hline
		
	\end{longtable}
	
	\section{Funkcje}
	
	\begin{longtable}{|p{16cm}|}
		\hline
		
		\textbf{Wymagania edukacyjne na ocenę dopuszczającą} \\ \hline
		
		\textbullet\ rozpoznaje przyporządkowania będące funkcjami w prostych przypadkach; \\
		
		\textbullet\ określa funkcję na różne sposoby: za pomocą grafu, tabeli, wykresu, opisu słownego i wzoru; \\
		
		\textbullet\ poprawnie stosuje pojęcia dziedziny, zbioru wartości, argumentu, miejsca zerowego, wartości i wykresu funkcji w prostych przypadkach; \\
		
		\textbullet\ odczytuje z nieskomplikowanego wykresu dziedzinę, zbiór wartości, miejsca zerowe oraz najmniejszą i największą wartość funkcji; \\
		
		\textbullet\ odczytuje z wykresu wartość funkcji dla danego argumentu oraz argument, dla którego funkcja przyjmuje daną wartość; \\
		
		\textbullet\ rozpoznaje wykresy funkcji rosnących, malejących i stałych; \\
		
		\textbullet\ oblicza wartość funkcji dla podanych argumentów na podstawie jej wzoru w prostych przypadkach; \\
		
		\textbullet\ wyznacza współrzędne punktu przecięcia wykresu funkcji danej wzorem z osią $OY$; \\
		
		\textbullet\ rozpoznaje wśród podanych wykresów wykresy funkcji $y=f(x-p)$, $y=f(x)+q$ oraz $y=f(x-p)+q$, gdy dany jest wykres funkcji $y=f(x)$; \\
		
		\textbullet\ wskazuje wielkości odwrotnie proporcjonalne. \\ \hline
		
		\textbf{Wymagania edukacyjne na ocenę dostateczną} \\ \hline
		
		\textbullet\ stosuje pojęcia dziedziny, zbioru wartości, argumentu, miejsca zerowego, wartości i wykresu funkcji; \\
		
		\textbullet\ na podstawie nieskomplikowanego wykresu funkcji określa argumenty, dla których funkcja przyjmuje wartości dodatnie, ujemne, niedodatnie i nieujemne; \\
		
		\textbullet\ określa na podstawie wykresu przedziały monotoniczności funkcji; \\
		
		\textbullet\ wyznacza dziedzinę funkcji określonej opisem słownym; \\
		
		\textbullet\ oblicza wartość funkcji dla różnych argumentów na podstawie jej wzoru; \\
		
		\textbullet\ oblicza argument odpowiadający podanej wartości funkcji w prostych przypadkach; \\
		
		\textbullet\ sprawdza algebraicznie, czy punkt o danych współrzędnych należy do wykresu funkcji danej wzorem; \\
		
		\textbullet\ wyznacza współrzędne punktów przecięcia wykresu funkcji danej wzorem z osią $OX$ w prostych przypadkach; \\
		
		\textbullet\ rysuje wykres funkcji danej wzorem w prostych przypadkach; \\
		
		\textbullet\ sporządza wykresy funkcji $y=f(x-p)$, $y=f(x)+q$ oraz $y=f(x-p)+q$ na podstawie wykresu funkcji $y=f(x)$; \\
		
		\textbullet\ stosuje funkcje i ich własności w prostych sytuacjach praktycznych; \\
		
		\textbullet\ wyznacza współczynnik proporcjonalności odwrotnej; \\
		
		\textbullet\ stosuje zależność między wielkościami odwrotnie proporcjonalnymi do rozwiązywania prostych zadań; \\
		
		\textbullet\ podaje wzór proporcjonalności odwrotnej, jeśli zna współrzędne punktu należącego do jej wykresu; \\
		
		\textbullet\ szkicuje wykres funkcji $f(x)=\frac{a}{x}$ dla danego $a>0$ i $x>0$. \\ \hline
		
		\textbf{Wymagania edukacyjne na ocenę dobrą} \\ \hline
		
		\textbullet\ rozpoznaje i opisuje zależności funkcyjne w sytuacjach praktycznych; \\
		
		\textbullet\ przedstawia daną funkcję na różne sposoby; \\
		
		\textbullet\ uzasadnia, dobierając odpowiednie argumenty, że funkcja nie jest monotoniczna; \\
		
		\textbullet\ odczytuje z wykresu funkcji rozwiązania równania $f(x)=m$ dla ustalonej wartości $m$; \\
		
		\textbullet\ odczytuje z wykresu funkcji zbiory rozwiązań nierówności $f(x)<m$, $f(x)>m$, $f(x)\leq m$ oraz $f(x)\geq m$ dla ustalonej wartości $m$; \\
		
		\textbullet\ szkicuje wykresy funkcji określonych różnymi wzorami w różnych przedziałach; \\
		
		\textbullet\ szkicuje wykresy funkcji, stosując przekształcenia wykresu w prostych przypadkach; \\
		
		\textbullet\ stosuje funkcje i ich własności, w tym proporcjonalność odwrotną, do rozwiązywania zadań praktycznych dotyczących drogi, prędkości i czasu w prostych przypadkach. \\ \hline
		
		\textbf{Wymagania edukacyjne na ocenę bardzo dobrą} \\ \hline
		
		\textbullet\ odczytuje z wykresów funkcji rozwiązania równań i nierówności typu $f(x)=g(x)$, $f(x)<g(x)$ oraz $f(x)>g(x)$; \\
		
		\textbullet\ szkicuje wykresy funkcji spełniających podane warunki w trudniejszych przypadkach oraz funkcji określonych różnymi wzorami w różnych przedziałach; \\
		
		\textbullet\ szkicuje wykresy funkcji, stosując przekształcenia wykresu w trudniejszych przypadkach; \\
		
		\textbullet\ stosuje funkcje i ich własności, w tym proporcjonalność odwrotną, do rozwiązywania zadań praktycznych dotyczących drogi, prędkości i czasu; \\
		
		\textbullet\ uzasadnia monotoniczność funkcji opisanej nieskomplikowanym wzorem, korzystając z definicji monotoniczności. \\ \hline
		
		\textbf{Wymagania edukacyjne na ocenę celującą} \\ \hline
		
		\textbullet\ rozwiązuje zadania o znacznym stopniu trudności dotyczące funkcji. \\ \hline
		
	\end{longtable}
	
	\section{Funkcja liniowa}
	
	\begin{longtable}{|p{16cm}|}
		\hline
		
		\textbf{Wymagania edukacyjne na ocenę dopuszczającą} \\ \hline
		
		\textbullet\ rozpoznaje funkcję liniową na podstawie jej wzoru lub wykresu; \\
		
		\textbullet\ rysuje wykres funkcji liniowej danej wzorem; \\
		
		\textbullet\ określa monotoniczność funkcji liniowej danej wzorem; \\
		
		\textbullet\ oblicza wartość funkcji liniowej dla danego argumentu; \\
		
		\textbullet\ wyznacza miejsce zerowe funkcji liniowej; \\
		
		\textbullet\ odczytuje z wykresu funkcji liniowej zbiór argumentów, dla których funkcja przyjmuje wartości dodatnie lub ujemne; \\
		
		\textbullet\ wyznacza współrzędne punktów przecięcia wykresu funkcji liniowej z osiami układu współrzędnych; \\
		
		\textbullet\ oblicza współczynnik kierunkowy prostej, jeśli dane są współrzędne dwóch punktów należących do tej prostej, a współrzędne tych punktów są liczbami wymiernymi; \\
		
		\textbullet\ rozpoznaje proste równoległe wśród prostych opisanych równaniami kierunkowymi; \\
		
		\textbullet\ rozpoznaje wielkości wprost proporcjonalne. \\ \hline
		
		\textbf{Wymagania edukacyjne na ocenę dostateczną} \\ \hline
		
		\textbullet\ interpretuje współczynniki występujące we wzorze funkcji liniowej; \\
		
		\textbullet\ oblicza argument, dla którego funkcja liniowa przyjmuje daną wartość; \\
		
		\textbullet\ wyznacza algebraicznie zbiory argumentów, dla których funkcja liniowa przyjmuje wartości dodatnie, ujemne, niedodatnie lub nieujemne; \\
		
		\textbullet\ wyznacza równanie prostej przechodzącej przez dwa dane punkty; \\
		
		\textbullet\ przekształca równanie ogólne prostej do postaci kierunkowej i odwrotnie; \\
		
		\textbullet\ sprawdza, czy trzy dane punkty są współliniowe; \\
		
		\textbullet\ rozpoznaje proste prostopadłe wśród prostych opisanych równaniami kierunkowymi; \\
		
		\textbullet\ wyznacza wzór funkcji liniowej, której wykres przechodzi przez dany punkt i jest równoległy do wykresu innej, danej funkcji liniowej; \\
		
		\textbullet\ wyznacza wzór funkcji liniowej, której wykres przechodzi przez dany punkt i jest prostopadły do wykresu innej, danej funkcji liniowej; \\
		
		\textbullet\ rozwiązuje układ równań metodą graficzną; \\
		
		\textbullet\ określa liczbę rozwiązań układu równań liniowych, korzystając z jego interpretacji geometrycznej; \\
		
		\textbullet\ opisuje równaniem zależność między wielkościami wprost proporcjonalnymi. \\ \hline
		
		\textbf{Wymagania edukacyjne na ocenę dobrą} \\ \hline
		
		\textbullet\ wyznacza wartości parametru, dla których funkcja liniowa jest rosnąca, malejąca lub stała; \\
		
		\textbullet\ określa wzajemne położenie prostych na podstawie ich równań; \\
		
		\textbullet\ wyznacza współrzędne wierzchołków wielokąta, gdy dane są równania prostych zawierających jego boki; \\
		
		\textbullet\ analizuje własności funkcji liniowej. \\ \hline
		
		\textbf{Wymagania edukacyjne na ocenę bardzo dobrą} \\ \hline
		
		\textbullet\ oblicza pole figury ograniczonej wykresami funkcji liniowych i osiami układu współrzędnych; \\
		
		\textbullet\ wyznacza wartości parametru, dla których dwie proste są równoległe; \\
		
		\textbullet\ stosuje warunki równoległości i prostopadłości prostych w dowodach dotyczących własności figur geometrycznych; \\
		
		\textbullet\ analizuje własności funkcji liniowej w zależności od wartości współczynników występujących w jej wzorze. \\ \hline
		
		\textbf{Wymagania edukacyjne na ocenę celującą} \\ \hline
		
		\textbullet\ wyprowadza wzór na współczynnik kierunkowy prostej przechodzącej przez dwa punkty; \\
		
		\textbullet\ udowadnia warunek prostopadłości prostych o danych równaniach kierunkowych; \\
		
		\textbullet\ rozwiązuje zadania o znacznym stopniu trudności dotyczące funkcji liniowej. \\ \hline
		
	\end{longtable}
	
	\section{Planimetria}
	
	\begin{longtable}{|p{16cm}|}
		\hline
		
		\textbf{Wymagania edukacyjne na ocenę dopuszczającą} \\ \hline
		
		\textbullet\ rozróżnia trójkąty ostrokątne, prostokątne i rozwartokątne; \\
		
		\textbullet\ stosuje twierdzenie o sumie miar kątów w trójkącie w prostych przypadkach; \\
		
		\textbullet\ uzasadnia przystawanie trójkątów, wykorzystując cechy przystawania, w prostych przypadkach; \\
		
		\textbullet\ zapisuje proporcje długości boków w trójkątach podobnych; \\
		
		\textbullet\ sprawdza, czy dane figury są podobne; \\
		
		\textbullet\ wskazuje w wielokątach odcinki, których długości są proporcjonalne. \\ \hline
		
		\textbf{Wymagania edukacyjne na ocenę dostateczną} \\ \hline
		
		\textbullet\ sprawdza, czy z trzech odcinków o danych długościach można zbudować trójkąt; \\
		
		\textbullet\ uzasadnia podobieństwo trójkątów, wykorzystując cechy podobieństwa, w prostych przypadkach; \\
		
		\textbullet\ wykorzystuje podobieństwo trójkątów do rozwiązywania prostych zadań; \\
		
		\textbullet\ stosuje podobieństwo wielokątów do obliczania długości boków; \\
		
		\textbullet\ stosuje twierdzenie o stosunku pól figur podobnych w prostych zadaniach; \\
		
		\textbullet\ rozwiązuje proste zadania, wykorzystując twierdzenie Talesa. \\ \hline
		
		\textbf{Wymagania edukacyjne na ocenę dobrą} \\ \hline
		
		\textbullet\ przeprowadza dowód twierdzenia o sumie miar kątów w trójkącie; \\
		
		\textbullet\ stosuje cechy przystawania trójkątów do rozwiązywania zadań geometrycznych; \\
		
		\textbullet\ wykorzystuje podobieństwo trójkątów do rozwiązywania problemów praktycznych i zadań geometrycznych; \\
		
		\textbullet\ rozwiązuje zadania dotyczące podobieństwa wielokątów; \\
		
		\textbullet\ rozwiązuje zadania, wykorzystując twierdzenie Talesa; \\
		
		\textbullet\ rozwiązuje zadania wymagające uzasadnienia i dowodzenia z wykorzystaniem przystawania trójkątów. \\ \hline
		
		\textbf{Wymagania edukacyjne na ocenę bardzo dobrą} \\ \hline
		
		\textbullet\ przeprowadza dowód twierdzenia o mierze kąta zewnętrznego trójkąta; \\
		
		\textbullet\ udowadnia, że symetralne boków trójkąta przecinają się w jednym punkcie; \\
		
		\textbullet\ udowadnia, że dwusieczne kątów trójkąta przecinają się w jednym punkcie; \\
		
		\textbullet\ stosuje cechy przystawania trójkątów do rozwiązywania trudniejszych zadań geometrycznych; \\
		
		\textbullet\ wykorzystuje podobieństwo trójkątów do rozwiązywania problemów praktycznych i trudniejszych zadań geometrycznych; \\
		
		\textbullet\ rozwiązuje zadania wymagające uzasadnienia i dowodzenia z wykorzystaniem podobieństwa trójkątów; \\
		
		\textbullet\ rozwiązuje zadania wymagające uzasadnienia i dowodzenia z zastosowaniem twierdzenia Talesa. \\ \hline
		
		\textbf{Wymagania edukacyjne na ocenę celującą} \\ \hline
		
		\textbullet\ przeprowadza dowód twierdzenia Talesa; \\
		
		\textbullet\ rozwiązuje zadania o znacznym stopniu trudności dotyczące przystawania i podobieństwa figur. \\ \hline
		
	\end{longtable}
	
\end{document}